\section{Step 22: Apply an accepted delta in place}
Restart at the payload beginning with pixel cursor zero. Read each command
again, advance the pixel cursor by its skip, and form the image address as the
base plus four times that cursor. For each changed pixel, pack the next triple
into a numeric word, load the old image word, XOR them, and store the result.
Advance source by three and destination by four. A zero changed count needs no
pixel loads or stores.

Fill \texttt{record\_delta}. Keep the illustrated pixel cursor in
\texttt{\$s7}. This second pass relies on the completed validation pass and
on the caller keeping the record immutable throughout the call. It requires no
rollback image. The existing XOR kernel explains the numeric operation, but the
record's packed triples first need conversion to native words; this scalar
construction performs that conversion inside the application loop.

\paragraph{Observe} Prepare \texttt{s22-1} and stop before
\texttt{record\_delta\_applied}. At \texttt{record\_delta\_applied}, \texttt{\$s7} is 4096.
The first image word changes from \texttt{0x000061c2} to
\texttt{0x000161c2}; the remaining 4095 words are unchanged. Publication
is still unfinished, so continuing returns development status four.

\Needspace{7\baselineskip}
\paragraph{Check} Why can this pass write directly into the working image?
\begin{itemize}
\item[A.] Every command has already been validated, and the record stays immutable.
\item[B.] XOR automatically undoes any later failure.
\item[C.] All delta records contain exactly one changed pixel.
\end{itemize}
